\chapter{Introduction and Scope}
\label{ch:introduction}

\section{The problem}

Pinhole glasses are a well-known, non-refractive way to improve a myopic
(short-sighted) eye's perceived sharpness: restricting the eye's entrance aperture
increases depth of focus, so a blurred retinal image sharpens even though the eye's
own optical power is unchanged. The IYPT-style problem statement motivating this
project's whole body of work is: \emph{investigate the optimal arrangement of pinholes
in an opaque film to produce sunglasses that reduce light intensity and help correct
short-sightedness.} That statement bundles two physically distinct design questions:

\begin{enumerate}
\item \textbf{Aperture size} --- for a single hole, what diameter minimizes retinal
  blur for a myopic eye, and how does a brightness (light-transmission) requirement
  change that answer?
\item \textbf{Aperture arrangement} --- if many holes are needed (to control total
  light throughput, field of view, or manufacturability), how should they be spaced
  and packed to avoid double-imaging (ghosting) between adjacent holes' sub-images?
\end{enumerate}

This report is scoped \textbf{entirely to question (1)}. The project's separate array/
packing/ghosting body of work (\texttt{array\_ghosting\_theory/}, \texttt{stage\_array/})
answers question (2) using a materially different theoretical apparatus (regime
classification, harmonic-series ghosting margins, a validated Barten 1999 contrast
model) and is documented in its own \texttt{RESULTS.md}; it is deliberately not
included here, per this report's own brief. Every model, sweep, and figure below
concerns a single circular aperture.

\section{What ``optimal'' means, and why this report never picks one number}

Four genuinely different objective functions appear across this project's history,
and they do not agree with one another:

\begin{itemize}
\item \textbf{Minimize retinal blur alone} (Chapters~\ref{ch:stage2}--\ref{ch:r1sweep}):
  treats the eye as a pure imaging system and asks which aperture produces the
  smallest encircled-energy blur radius at the retina. This objective, taken alone,
  favors apertures far smaller than what is clinically observed to work.
\item \textbf{Minimize blur subject to a brightness floor}
  (Chapter~\ref{ch:stage3a}): the same blur objective, but only over apertures that
  admit at least some minimum fraction of full-pupil illuminance. This single added
  constraint moves the optimum by roughly $2$--$3\times$ into the clinically observed
  bracket.
\item \textbf{Maximize predicted visual acuity} (Chapters~\ref{ch:acuity_models}--\ref{ch:stage4}):
  folds in contrast sensitivity, photometry, and pupil dynamics to predict a clinical
  logMAR score directly, rather than a physical blur radius.
\item \textbf{Match observed clinical acuity data} (Chapters~\ref{ch:r7acuity}--\ref{ch:validation}):
  the strictest bar --- does any of the above models, however built, actually predict
  the two real clinical datasets available? This report's honest answer, reached
  independently by two structurally different models, is: not yet, on the data
  currently available.
\end{itemize}

Because these four objectives disagree, and because the clinical-validation bar is not
currently met by any model in this project, no chapter in this report issues an
aperture recommendation. Each chapter reports precisely what its own model predicts,
under its own explicitly stated assumptions; Chapter~\ref{ch:synthesis} lays every
model's answer side by side without collapsing them into one.

\section{How this report is organized}

Part I (Chapters~\ref{ch:optical_system}--\ref{ch:wave_optics}) establishes the shared
optical and myopia model every later chapter builds on. Part II
(Chapters~\ref{ch:stage1}--\ref{ch:r1sweep}) is the core wave-optics blur-minimization
pipeline, including a full account of a real, confirmed bug in that pipeline and how
fixing it changed the headline result. Part III
(Chapters~\ref{ch:stage3a}--\ref{ch:singleilluminate}) covers brightness and
illumination constraints. Part IV (Chapters~\ref{ch:mtf_strehl}--\ref{ch:stage3_integrated})
covers the human-vision layer: MTF, contrast sensitivity, and acuity prediction. Part V
(Chapters~\ref{ch:classical_model}--\ref{ch:validation}) covers model validation,
including a second, independent acuity architecture and the clinical-data comparison.
Part VI (Chapters~\ref{ch:percentile}--\ref{ch:limitations}) is the synthesis: every
model's optimal-aperture answer laid out together, the corrected-vs-deprecated optics
comparison, and an explicit, itemized list of this project's own known limitations.

\section{Data and reproducibility}

Every sweep in this report was generated by directly calling this project's own
already-implemented, already-tested compute functions (never a re-derivation) --- the
same thin-wrapper functions that back this project's live interactive dashboard
(\texttt{master\_dashboard/wrappers.py}), so the numbers here are the same numbers a
reader could reproduce by running that dashboard themselves. Sweeps use a common
$15\times 8$ grid throughout: aperture $a_{hole}\in\{75,100,150,200,300,400,500,600,
700,750,1000,1250,1500,1750,2000\}\um$ and myopia $M\in\{0,1,2,3,4,5,6,8\}\,$D, at
wavelength $550\,$nm, pinhole--cornea gap $12\mm$, and natural pupil diameter
$4000\um$ unless a chapter states otherwise. Appendix~\ref{app:software} documents the
exact software architecture and the specific safeguard (an isolated, self-restoring
Python import context) used to guarantee the deprecated-optics comparison in
Chapter~\ref{ch:deprecated_vs_corrected} never contaminates any other chapter's numbers.
